\documentclass[11pt]{article}
\usepackage[T1]{fontenc}
\usepackage[utf8]{inputenc}
\usepackage{lmodern}
\usepackage[a4paper,margin=1in]{geometry}
\usepackage{microtype}
\usepackage{amsmath,amssymb,amsthm,mathtools}
\usepackage{booktabs,array,longtable}
\usepackage[dvipsnames]{xcolor}
\usepackage[colorlinks=true,linkcolor=MidnightBlue,urlcolor=MidnightBlue,citecolor=MidnightBlue]{hyperref}
\usepackage{enumitem}
\usepackage{setspace}
\onehalfspacing
\setlist[itemize]{topsep=4pt,itemsep=2pt,leftmargin=1.5em}
\setlist[enumerate]{topsep=4pt,itemsep=2pt,leftmargin=1.8em}

\newcommand{\DR}{\mathsf{DR}}
\newcommand{\PO}{\mathsf{PO}}
\newcommand{\PP}{\mathsf{PP}}
\newcommand{\PPI}{\mathsf{PPI}}
\newcommand{\EQ}{\mathsf{EQ}}
\newcommand{\Base}{\{\DR,\PO,\EQ,\PP,\PPI\}}
\newcommand{\cl}{\operatorname{cl}}
\newcommand{\tp}{\operatorname{tp}}
\newcommand{\st}{\operatorname{st}}
\newcommand{\comp}{\operatorname{comp}}
\newcommand{\parQ}{\mathsf{par}_Q}
\newcommand{\eqQ}{\equiv_{\EQ,Q}}
\newcommand{\eqcong}{\equiv_{\EQ}}
\newcommand{\blk}{\equiv_{\mathrm{blk}}}
\newcommand{\reach}{\rightsquigarrow}
\newcommand{\Mate}{\operatorname{Mate}}
\newcommand{\ALCIRCC}{\mathcal{ALCI}_{\mathsf{RCC5}}}

\title{\textbf{Formal Review of the Latest v2.1 Responses}\\[0.35em]
\large Remaining gaps in the proposed split-forest decidability proof for $\ALCIRCC$}
\author{Review memorandum}
\date{May 2026}

\begin{document}
\maketitle

\begin{abstract}
This memorandum reviews the latest submitted responses,
\texttt{completeness\_extraction\_ALCIRCC5\_v2(1).tex} and
\texttt{sibling\_interface\_descriptors\_ALCIRCC5\_v2(1).tex}, together
with their PDFs.  The update is a real improvement: the density/Hasse
objection is now treated in the right way via witness-generated submodels
and generated cover presentations; $\EQ$-modalities are normalized under
strong equality; and the documents try to separate blocking from semantic
$\EQ$.  However, the proof still does not establish decidability as written.
Several of the claimed v2.1 repairs introduce new formal problems.  The most
serious are: a contradiction between profile-level equality ports, self-loop
blocking, vertical separation, and mosaic reflexivity; an invalid finite-index
argument for equality ports; a still-unresolved multiple-upward-$\PP$-witness
problem; an incorrect assumption that equality mates have a unique semantic
parent; a canonical-unfolding/projection definition that is not compatible
with the rootless containment orientation; and an overstrong finite-mosaic
arity claim.
\end{abstract}

\section{Executive assessment}

The latest response should be viewed as a useful draft of a repair, not as a
closed proof.  It correctly moves the proof away from arbitrary semantic Hasse
diagrams and toward the intended model class: sparse witness-generated abstract
models with a generated containment presentation.  This substantially improves
the earlier proof architecture.

Nevertheless, I would not yet rely on the latest v2.1 documents as a proof of
\(\ALCIRCC\) decidability.  The central issue is not the old density objection;
that issue is now essentially resolved.  The remaining issues are internal to
the finite certificate itself, especially the interaction between:
\[
\text{blocking/profile repetition},\qquad
\text{semantic }\EQ\text{ quotienting},\qquad
\text{multiple generated superpart witnesses}.
\]
The current equality-port repair is inconsistent with the self-loop certificate
needed for the motivating formula
\[
C_\uparrow = \exists\PP.\top \sqcap \forall\PP.\exists\PP.\top.
\]
Also, the current parent-function treatment cannot yet support formulas whose
satisfaction requires several incomparable proper-superpart witnesses for the
same semantic element.

\section{What has improved}

The following changes are positive and should be retained.

\subsection{Witness-generated submodels}

The witness-generated submodel lemma in
\texttt{completeness\_extraction\_ALCIRCC5\_v2(1).tex}, lines 430--512, is the
right conservativity step.  It reduces arbitrary satisfying abstract models to
sparse models generated by selected existential witnesses.  The closure is now
closed under NNF complement, so the converse direction for universal formulas
uses the selected witness for \(\exists R.\overline F\) correctly.

This largely removes the old density/Hasse concern.  The proof no longer needs
to assume that the original semantic \(\PP/\PPI\) relation has immediate covers.

\subsection{Generated covers rather than ambient Hasse covers}

The generated-cover reinterpretation in lines 520--570 is also the right
conceptual move.  A tree edge should be a generated witness-cover edge relative
to the selected presentation, not a Hasse edge of an arbitrary ambient
\(\PP\)-order.

\subsection{Elimination of $\EQ$ modalities under strong equality}

The normalization lemma in lines 302--358 correctly observes that, under strong
\(\EQ\),
\[
\exists\EQ.D \equiv D,\qquad \forall\EQ.D \equiv D.
\]
This is a good simplification.  After this preprocessing, the modal roles that
need witness handling are only
\(\DR,\PO,\PP,\PPI\).

\subsection{Recognition that blocking is not semantic equality}

The v2.1 text explicitly tries to distinguish finite blocking from semantic
\(\EQ\), for example around lines 752--767 and 1196--1381.  This distinction is
essential.  Unfortunately, the particular equality-port implementation is still
not coherent, as explained next.

\section{Critical remaining gaps}

\subsection{G1. Profile-level equality ports contradict self-loop blocking}

This is the most serious defect in the current v2.1 repair.

The finite quotient explicitly allows a parent self-loop
\(\parQ(\sigma)=\sigma\) to represent an infinite upward \(\PP\)-chain; see
lines 752--767.  The reachability definition then says that a self-loop implies
\(\sigma\reach_{\PP}\sigma\); see lines 1718--1724.

The new vertical-separation axiom says that if \(\sigma\reach_{\PP}\sigma'\) or
\(\sigma\reach_{\PPI}\sigma'\) strictly, then no profile-port pair
\((\sigma,p)\) may be \(\eqQ\)-related to \((\sigma',q)\); see lines
1274--1285.  Therefore, for a self-loop it entails
\[
(\sigma,p)\not\eqQ(\sigma,p)
\]
for every port \(p\).  The proof of Lemma 10.5 explicitly uses this conclusion
at lines 1326--1332.

But mosaics also require reflexive equality.  Axiom (M1) says
\(\mu(v,v)=\EQ\) for every vertex; see lines 843--849.  The revised axiom (M7)
says
\[
\mu(v,w)=\EQ \quad\Longleftrightarrow\quad v\eqQ^\mu w
\]
for all vertices \(v,w\); see lines 1252--1260.  Taking \(v=w\), M1 and M7
force
\[
(\st(v),\pi(v))\eqQ(\st(v),\pi(v)).
\]
For a vertex of self-loop profile \(\sigma\), this is exactly the relation
that vertical separation forbids.

Thus the one-state certificate for \(C_\uparrow\) is made invalid by the very
repair intended to support it.  The problem is structural: a single finite
relation on profile-port pairs cannot distinguish ``the same occurrence'' from
``the next lap of the blocking cycle'' when both carry the same profile and the
same port tuple.

\paragraph{Required repair.}
The equality mechanism must be occurrence-sensitive or context-sensitive.  One
possible direction is:
\begin{itemize}
\item treat \(v=w\) in mosaics as identity, not as a global \(\eqQ\) lookup;
\item use finite equality variables local to a bounded mosaic/context, not a
global relation on reusable profile-port pairs; and
\item state explicitly that equality links are between distinct displayed
occurrences in a finite context, while blocking repetition never induces such a
link across cycle laps.
\end{itemize}
As written, the current global profile-port relation \(\eqQ\) cannot serve both
roles.

\subsection{G2. The finite equality-port bound is not proved and appears false}

Lines 1348--1355 say that equality ports are drawn from a finite alphabet whose
size is bounded by the number of distinct semantic elements that can appear in
a witness-generated submodel cell, and hence by a computable function of
\(|\cl(C_0)|\).  Lines 1383--1408 then define the extracted congruence by using
the origin element itself as the canonical identity index.

This does not yield a finite alphabet.  A witness-generated model may contain
infinitely many distinct semantic elements of the same rank-\(d\) profile.  In
strong \(\EQ\) semantics, each distinct element is its own equality class.  If
one port is needed per origin to avoid false equality, the port alphabet is
infinite.  If finitely many ports are reused for distinct origins, then M7 may
force false \(\EQ\) labels when two reused ports appear together in a mosaic.

The sentence in lines 1405--1408 only bounds the number of ports contributed by
one semantic element; it does not bound the number of semantic elements.  Hence
Lemma 10.8 does not establish the finite-index property.

\paragraph{Required repair.}
Do not try to encode global semantic identity by finite reusable profile ports.
Instead, equality should be stored as bounded local matching data for the finite
interfaces that actually appear in a certificate.  If global equality classes
are needed, a separate finite abstraction theorem must prove that only boundedly
many equality representatives per profile are ever relevant.

\subsection{G3. Multiple upward $\PP$ witnesses are still not handled}

The containment orientation makes \(\exists\PP.D\) point upward to a proper
superpart.  A single occurrence has only one parent.  However, a semantic element
may need several incomparable proper-superpart witnesses.

The split-tree construction in lines 595--604 assigns a single parent to an
occurrence.  For each \(\exists\PP.D\), it either uses the existing parent or
sets \(\mathsf{par}(u)=u^\uparrow\).  This does not define what happens when
one occurrence has two incompatible upward witnesses.  The later splitting
clause, lines 620--624, only says that if a single element \(b\) must appear
under multiple parents, then one should introduce one occurrence per containment
branch.  The problematic case is the dual one: the same lower element \(a\)
may need to appear under several selected superparts.

A representative stress formula is:
\[
\begin{aligned}
C_{AB}:={}&\exists\PP.A \sqcap \exists\PP.B \\
&{}\sqcap \forall\PP\bigl(A\rightarrow(\forall\PP.\neg B\sqcap\forall\PPI.\neg B)\bigr) \\
&{}\sqcap \forall\PP\bigl(B\rightarrow(\forall\PP.\neg A\sqcap\forall\PPI.\neg A)\bigr).
\end{aligned}
\]
This is satisfiable by an element \(x\) with two incomparable proper superparts
\(y,z\), with \(A(y)\), \(B(z)\), and \(y\PO z\).  A single parent chain above
one occurrence of \(x\) cannot contain both witnesses, because the formula
forbids the \(A\)- and \(B\)-witnesses from being comparable.  The intended
split-copy idea can handle this only by creating two occurrences of \(x\), one
under \(y\) and one under \(z\), and later quotienting them as semantic equality
mates.

The current v2.1 text has not formally defined this construction.  In fact,
Lemma 12.3 goes in the opposite direction: lines 1968--1971 say that equality
mates have either no parent or the same unique semantic parent.  But the whole
reason to split is that equality mates of the same semantic element may need
different generated parents.

\paragraph{Required repair.}
The parent function should be a function on occurrences, not on semantic
equality classes.  Equality mates must be allowed to have different generated
parents.  Therefore \(\parQ\) should not be required to descend to mate classes.
Eventuality discharge should quantify over equality-mate occurrences or over a
finite set of parent contexts, without imposing parent congruence.

\subsection{G4. The canonical unfolding/projection is not compatible with the rootless orientation}

The canonical unfolding definition in lines 1865--1875 says that \(U(Q)\) is a
rooted tree of finite walks in \((S,\parQ)\) starting at the evaluation profile
\(\sigma_\ast\), with parent map given by one-step predecessor.  The projection
\(\pi:T\to U(Q)\) in lines 1878--1888 uses the depth of an occurrence relative
to a root occurrence \(u_\ast\).

This is not compatible with the containment-oriented construction.  The
evaluation occurrence may have generated parents above it; indeed this is
necessary for \(\exists\PP\)-requirements.  After upward expansion, \(u_\ast\)
is not a root of the containment tree.  Ancestors of \(u_\ast\) do not have a
finite depth below \(u_\ast\), and the displayed formula for \(\pi(u)\) is not
defined for them.

There is also an orientation problem: if \(\parQ(\sigma_i)=\sigma_{i+1}\), then
the walk \((\sigma_\ast,\sigma_1,\ldots,\sigma_m)\) moves upward along the
containment parent function.  But the parent map of the rooted walk tree is
truncation to the predecessor, which reverses the intended semantic parent
relation.  This may be a notational repair, but as written the projection is
not a faithful unfolding map.

The surjectivity claim in Lemma 12.2 is also not justified.  Given an arbitrary
finite walk in \((S,\parQ)\), choosing any occurrence with the endpoint profile
does not ensure that its actual parent chain realizes the specified preceding
profiles.

\paragraph{Required repair.}
Use a pointed bi-oriented unfolding: nodes should be finite paths extending both
above and below the evaluation occurrence, or else represent upward trunks as
one-hole contexts/lassos rather than rooted walks from \(\sigma_\ast\).  The
projection must be defined on occurrence paths, not merely on endpoint profiles.

\subsection{G5. The side-context witness lemma remains circular}

The side-context witness lemma in lines 1007--1185 is an improvement in form,
but its proof still assumes the placement it is supposed to prove.

The witness-menu extraction in lines 961--1003 says that a \(\DR/\PO\) witness
occurrence exists, ``possibly as a sibling occurrence under a common ancestor.''
The side-context lemma then proves existence of such a \(v\) by citing
Construction 5.1, step 4.  But step 4, lines 614--618, merely says to record an
outgoing witness edge to a node mapped to \(b\), ``typically a sibling of \(u\)
under a common ancestor.''  It does not construct the occurrence, the common
ancestor, or the corresponding containment placement.

The use of a lowest common ancestor also requires care.  In a rootless or
upward-regular containment presentation, and especially after split copies,
there need not be a unique LCA in the naive rooted-tree sense.  The lemma should
not use an LCA until the underlying generated presentation has been formally
shown to be an actual tree with the relevant nodes in the same component and
with a well-defined meeting point.

Finally, a five-vertex side context checks only a bounded immediate interface.
It does not by itself prove that all deeper universal restrictions and all
composition constraints involving descendants of \(u\) or \(v\) remain coherent.
The text says these are handled by support closure, but the support-closure
argument is itself still incomplete; see G6.

\paragraph{Required repair.}
State and prove a genuine placement theorem: for every selected \(\DR/\PO\)
witness, the split presentation can be extended by a finite side context such
that all induced vertical, sibling, and equality constraints are preserved.  The
proof must construct the context rather than cite the schematic step 4.

\subsection{G6. The finite mosaic arity theorem is overclaimed}

Theorem 9.3, lines 1619--1660, claims that arity \(5\) suffices for all mosaics,
using three ingredients: local arity of axioms (M1)--(M7), local arity of
support-closure operations (SC1)--(SC4), and the RCC5 patchwork property.

This does not yet prove the needed result.  Checking that each axiom is binary
or ternary is not the same as proving that descriptor membership, closure under
iterated extensions, and global realizability are determined by all submosaics
of size at most \(5\).  In particular:
\begin{itemize}
\item SC4 can be iterated along containment chains.  The proof says that adding
one generated child raises arity by one and stays within arity \(5\), but an
iterated construction may couple several extensions through shared endpoints.
\item The patchwork property is an amalgamation principle for compatible
networks; it does not automatically imply a finite forbidden-subnetwork theorem
of arity \(5\) for the enriched descriptor language with statuses, equality
ports, witness menus, and support closure.
\item The theorem is used to justify finite enumeration of support-closed
mosaics, but the closure object is defined semantically as containing joint
networks occurring in a model.  The equivalence between semantic occurrence and
bounded syntactic closure is not proved.
\end{itemize}

\paragraph{Required repair.}
Either give a precise finite closure algorithm and prove by induction that it
covers every finite prefix network needed in the soundness construction, or
replace the constant \(5\) by a computable bound derived from the finite state
space and prove a finite-model/compactness bound for descriptors.  The current
patchwork citation is not sufficient.

\subsection{G7. Simultaneous realization still relies on the hidden source model}

Lemma 12.4, lines 1978--2099, tries to prove simultaneous realization by taking
\(U^\ast(Q)=T\), where \(T\) is the split tree extracted from the original model.
For the model-to-quotient completeness direction, this can be useful as a
witness that the extracted data were realizable.  But the decision procedure
needs a finite validity condition that guarantees realization for an arbitrary
enumerated finite certificate, without access to a hidden model \(T\).

As written, the proof of simultaneous realization repeatedly appeals to the
source model \(\mathcal I_w\), the original witness selector, and the origin map.
For example, SR1 uses the existence of a witness in \(\mathcal I_w\) and then
asserts that the split-tree construction installs a finite generated cover chain
visiting it.  This proves at most that the particular extracted tree worked,
assuming the earlier construction was correct.  It does not provide a decidable
certificate-side criterion independent of the model.

This matters because the main theorem later combines completeness with
soundness by enumerating finite quotients.  The finite validity check must be
strong enough that every accepted quotient has a coherent unfolding.  The
current argument does not yet bridge that gap, especially given G1--G6.

\paragraph{Required repair.}
Define validity purely on finite data and prove a soundness theorem from that
finite validity condition to an unfolded model.  If reachability is the chosen
replacement for lasso/Buechi acceptance, then the proof must show that finite
reachability plus the finite mosaic/equality conditions entail one simultaneous
unfolding for all eventualities, not merely for extracted quotients.

\subsection{G8. The sibling-interface note is out of sync with v2.1}

The sibling-interface document remains useful as a high-level compatibility
note, but it does not yet match the latest v2.1 completeness manuscript.

First, lines 459--465 state a simple reachability-discharge condition from
\(\sigma\) alone.  The v2.1 completeness manuscript changed this to
mate-class-aware reachability.  The sibling note therefore still states the old
v2 condition, not the claimed v2.1 condition.

Second, lines 265--268 allow a literal parent self-loop in the generated
split-tree presentation.  A literal occurrence-level self-loop would violate
irreflexivity of \(\PP/\PPI\).  Self-loops should exist only in the finite
quotient/blocking graph, whose unfolding creates fresh occurrences at each lap.
The sibling note says this later, but the definition itself remains misleading.

Third, the sibling note claims the v1 soundness chain is unchanged, while the
v2.1 completeness paper substantially changes equality handling.  If typed
\(\EQ\), M7, equality ports, and vertical separation are part of quotient
validity, then the soundness chain must be restated against those actual v2.1
validity conditions.

\paragraph{Required repair.}
Update the sibling-interface note to v2.1 terminology.  In particular, remove
literal occurrence-level self-loops, replace simple reachability by the final
correct equality-aware eventuality condition, and restate the soundness theorem
for the actual equality mechanism used by the completeness proof.

\section{Additional technical concerns}

\subsection{The open-cross-edge lemma is still too compressed}

The open-cross-edge lemma in lines 808--825 of the completeness manuscript is a
key structural claim: non-core descendants of sibling subtrees can have only
\(\DR\) or \(\PO\) cross-relations.  The proof is a short appeal to the old
open-cross-edge theorem.  But after generated covers, split copies, equality
ports, and mate-aware reachability are introduced, the notion of core must be
rechecked.  In particular, if equality mates may appear under different parents,
then whether a node is ``visible to both sides'' can depend on the chosen
occurrence context.  The old lemma may still be true, but it needs to be
restated and proved under the new occurrence/equality semantics.

\subsection{The RCC5 composition table is no longer displayed}

The latest completeness manuscript now cites a standard RCC5 composition table
instead of printing the earlier incorrect one.  This avoids the previous error,
but it also makes the paper less self-contained.  Since the proof relies on
composition at several critical points, the final manuscript should include the
correct atomic RCC5 table, or at least cite a fixed external source and state
which convention is used.

\subsection{The phrase ``all six obligations discharged'' is too strong}

The v2.1 changelog says that all six review obligations are discharged.  Given
G1--G8, that statement is not yet supportable.  A safer statement would be:

\begin{quote}
The present revision proposes repairs for the six obligations identified in the
v2 review.  Several repairs, especially the explicit equality-port mechanism,
the multiple-parent split construction, and the finite mosaic arity bound, still
require further formal verification.
\end{quote}

\section{Recommended path to a correct proof}

The current direction is promising, but the certificate needs a different
formalization of equality and upward witnesses.  I recommend the following
repair sequence.

\subsection{Step 1: Separate three relations completely}

The proof should maintain three distinct notions:
\[
\begin{array}{rcl}
\text{blocking repetition} &:& \text{same finite profile in the quotient},\\
\text{occurrence identity} &:& \text{the same vertex in a finite unfolding/context},\\
\text{semantic }\EQ &:& \text{a declared typed congruence between displayed occurrences}.
\end{array}
\]
A reusable profile-port pair cannot by itself encode occurrence identity.  M7
should not require global \(\eqQ\)-reflexivity on profile-port pairs when those
pairs are reused across blocking laps.

\subsection{Step 2: Make split copies explicit for upward witnesses}

For each occurrence and each upward \(\PP\)-eventuality, the construction should
be allowed to create a fresh occurrence-copy of the lower element under the
chosen superpart.  These copies may later be semantically quotiented, but their
parents need not agree.  Therefore parent preservation must not be an axiom of
semantic \(\EQ\).

A finite certificate can store this as a finite set of parent contexts rather
than a single parent function per semantic equality class.

\subsection{Step 3: Redefine the unfolding as a generated occurrence grammar}

Instead of rooted walks from \(\sigma_\ast\), use a finite occurrence grammar or
context graph whose generated objects are occurrence paths around the evaluation
point.  Blocking cycles in this grammar create fresh occurrences.  Equality
links are local edges between generated occurrences, subject to typed-congruence
checks, not profile-level identities.

\subsection{Step 4: Prove finite descriptor locality independently}

The support-closed mosaic construction should be given as an explicit finite
closure algorithm.  The proof should show, by induction on finite prefixes, that
every required finite prefix network is represented by the finite descriptor
language.  If a constant arity bound is claimed, prove it directly; otherwise
use a larger but computable bound.

\subsection{Step 5: Re-sync the sibling-interface and completeness papers}

Once the equality and parent-context formalism is fixed, the sibling-interface
note should be rewritten to match the final definitions exactly.  It should not
refer to v1 soundness as unchanged unless every altered v2.1 validity condition
has been checked against the v1 soundness proof.

\section{Conclusion}

The latest response is materially better than the earlier v2 draft.  It
correctly removes the density/Hasse obstruction and adopts the intended sparse,
witness-generated model class.  However, the decidability proof still does not
stand as written.

The decisive remaining problem is the finite certificate's equality and parent
structure.  The new \(\eqQ\) mechanism simultaneously needs to be reflexive for
mosaic self-edges and non-reflexive across self-loop blocking laps; a global
relation on profile-port pairs cannot do both.  Moreover, equality mates must be
allowed to have different generated parents in order to realize multiple
incomparable upward \(\PP\)-witnesses, whereas the current proof assumes a
unique semantic parent.

Thus the current v2.1 manuscripts should be treated as a useful repair proposal,
not as a complete proof.  The proof can likely be repaired, but it needs a
cleaner occurrence-level equality formalism and an explicit parent-context
construction for split copies before the finite quotient theorem can be
accepted.

\end{document}
